\section{Discusión}

En esta sección se interpretan los resultados obtenidos, se examina su relación con el estado del arte
y se analizan sus implicaciones para la adhesión fibra–matriz evaluada mediante el método
\textit{Fiber-Bundle Pull-Out} (FBPO). La discusión se organiza en torno a la interpretación de los
hallazgos experimentales, su comparación con la literatura, las limitaciones del estudio y las posibles
líneas de investigación futura.

\subsection{Interpretación de los resultados}

En esta subsección se analizan los resultados presentados en el capítulo anterior, considerando su
coherencia interna, los patrones observados y la respuesta de cada sistema fibra–matriz. Se evalúa la
magnitud relativa de la IFSS aparente, las diferencias entre sistemas y cualquier comportamiento
atípico registrado durante los ensayos.

\subsection{Relación con el estado del arte}

Aquí se comparan los resultados obtenidos con los rangos típicos reportados en la literatura para
ensayos FBPO y para métodos micromecánicos afines. Se destacan similitudes y discrepancias respecto de
estudios previos, analizando posibles causas de las diferencias observadas, tales como variaciones de
\textit{sizing}, compactación del haz, procedimiento de curado o condiciones experimentales.

\subsection{Limitaciones del estudio}

Esta subsección identifica y describe las principales limitaciones que pueden influir en la
interpretación de los resultados, tales como variabilidad experimental, número de muestras válidas,
restricciones de montaje, sensibilidad a la sujeción y precisión metrológica en la medición de la
longitud embebida y del diámetro equivalente del haz. Se evalúa el impacto potencial de estas
limitaciones y se discute cómo podrían abordarse en trabajos futuros.

\subsection{Implicancias y proyecciones}

En esta parte se discuten las implicancias prácticas y científicas de los resultados, considerando su
relevancia para el diseño de materiales compuestos, la selección de \textit{sizing} y la comparación entre
proveedores de fibra. Asimismo, se plantean recomendaciones para investigaciones futuras, tales como
optimización del protocolo experimental, ampliación del tamaño muestral, incorporación de técnicas
complementarias de caracterización superficial o implementación de modelos numéricos para describir el
comportamiento interfacial.




